De term bestandssysteem\index{bestandssysteem} of in het Engels file system\index{file system} (filesystem) is een complexe term. Het kan naar verschillende dingen verwijzen. Normaal gesproken wordt de term bestandssysteem gebruikt voor de manier waarop we bestanden opslaan op disk. Ook de hele boom aan directories noemen we een file system hoewel dat in een GNU/Linux systeem dus meerdere echte en virtuele disks kan beslaan.

Een bestandssysteem is een manier om data te organiseren op een opslag apparaat. Verschillende besturingssystemen hebben in de loop van de tijd verschillende bestandssystemen gebruikt. Microsoft is met DOS ooit begonnen met het \textit{msdos} bestandssysteem. Bestandsnamen bestonden daarbij uit 8 karakters, een punt en een 3 karakter extensie (filename.txt). Op GNU/Linux wordt er veel gebruik gemaakt van het \textit{ext} filesystem. De huidige versie is versie 4 vandaar de naam \textit{ext4}.

